\section{Introduction}\label{sec:intro}

When Eringen's integral nonlocal elasticity \cite{Eringen1972Linear,Eringen2002} is coupled with classical or generalized heat conduction to model nanoscale thermoelastic waves and resonators, the nonlocal attenuation operator is routinely replaced by the second-order Helmholtz differential operator $1 - \ell^2\partial_x^2$ (where $\ell = e_0 a$ is the internal characteristic length). For a homogeneous isotropic rod occupying a domain $\Omega \subseteq \Rbb$ with displacement $u(x,t)$, linearized strain $\varepsilon = \partial_x u$, and temperature increment $\theta = T - T_0$ ($|\theta|/T_0 \ll 1$), the balances of linear momentum and linearized entropy are
\begin{equation}\label{eq:balance_laws}
  \partial_x \sigma(x,t) + b(x,t) = \rho \ddot{u}(x,t), \qquad T_0 \rho \dot{\eta}(x,t) = -\partial_x q(x,t) + r(x,t),
\end{equation}
where $\sigma$ is the Cauchy stress, $\eta$ is the specific entropy density, $q$ is the heat flux, and $\rho, E, \beta, c_E, \kappa > 0$ denote the mass density, Young's modulus, thermal stress modulus, specific heat at constant strain, and thermal conductivity.

Although Eringen \cite{Eringen1974} derived the thermodynamically consistent integral equations of linear nonlocal thermoelasticity in 1974, subsequent applied boundary-value and wave-propagation studies have diverged into conflicting differential shortcuts regarding (i)~whether the nonlocal operator acts on the thermal stress $-\beta\theta$ without acting on the dilational heating $T_0\beta\dot{\varepsilon}$, (ii)~whether the purely nonlocal ($\gamma = 0$) differential operator is applied to all constitutive equations simultaneously, and (iii)~how the integral convolution is closed on bounded domains $\Omega = [0,L]$. This clarification note synthesizes the thermodynamic, spectral, and boundary-value conditions required to make two-phase nonlocal thermoelasticity well-posed and self-consistent.

\subsection{Prior Art and Theoretical Foundations}\label{subsec:prior_art}

Every mathematical ingredient used in this note rests on established prior art in nonlocal continuum mechanics and thermodynamics, which we credit explicitly:

\begin{enumerate}[leftmargin=*,align=left]
  \item \textbf{Edelen--Laws nonlocality residuals and Eringen's (1974) nonlocal thermoelasticity \cite{Eringen1972Linear,Eringen1974,Eringen2002}:} Building on the nonlocal thermodynamics of Edelen and Laws (1971), in which pointwise balance laws include zero-mean localization residuals ($\int_\Omega \hat{e}\dd x = 0$, $\int_\Omega \hat{h}\dd x = 0$), Eringen \cite{Eringen1974} (Section~3, Eqs.~(3.11)--(3.15)) proved that when the nonlocal Helmholtz free energy contains a bilinear thermomechanical coupling with a symmetric kernel $\alpha(|x-\xi|)$, the stress and entropy state laws must share the \emph{identical} integral operator:
  \begin{equation}\label{eq:eringen_1974_laws}
    t_{kl}(\bm{x},t) = \int_\Omega \Big[\lambda'(\lvert\bm{x}-\bm{\xi}\rvert)\varepsilon_{rr}(\bm{\xi},t)\delta_{kl} + 2\mu'(\lvert\bm{x}-\bm{\xi}\rvert)\varepsilon_{kl}(\bm{\xi},t) - \beta'(\lvert\bm{x}-\bm{\xi}\rvert)\theta(\bm{\xi},t)\Big]\dd V(\bm{\xi}),
  \end{equation}
  paired with $\rho\eta(\bm{x},t) = \int_\Omega \big[\beta'(\lvert\bm{x}-\bm{\xi}\rvert)\varepsilon_{kk}(\bm{\xi},t) + \frac{\rho c'}{T_0}(\lvert\bm{x}-\bm{\xi}\rvert)\theta(\bm{\xi},t)\big]\dd V(\bm{\xi})$. Thus, if the thermal stress is nonlocal with kernel $\beta'(|x-\xi|)$, the dilational entropy must carry the exact same kernel.

  \item \textbf{Polizzotto's variational thermodynamics and Borino--Failla--Parrinello (BFP) symmetric kernel normalization \cite{BazantJirasek2002,DiLuzio2005,MarottiDeSciarra2009,BazantLeHoover2010}:} Polizzotto (2001) reformulated Edelen--Laws nonlocal thermodynamics in Hilbert space via self-adjoint integral operators and variational principles. On a bounded domain $\Omega$, the unnormalized kernel $K_\ell^{\mathrm{unnorm}}(x,\xi) = \frac{1}{2\ell}\mathrm{e}^{-|x-\xi|/\ell}$ satisfies $\int_0^L K_\ell^{\mathrm{unnorm}}(x,\xi)\dd\xi = 1 - \vartheta_\ell(x) < 1$ with boundary weight defect $\vartheta_\ell(x) := \frac{1}{2}(\mathrm{e}^{-x/\ell} + \mathrm{e}^{-(L-x)/\ell})$, violating constant-field preservation near $\partial\Omega$, whereas naive row scaling $K_\ell(x,\xi)/\int_\Omega K_\ell(x,z)\dd z$ destroys symmetry and self-adjointness \cite{BazantJirasek2002,DiLuzio2005}. Borino, Failla, and Parrinello (2003) restored both symmetry and exact partition of unity ($\int_\Omega W_\ell^{\mathrm{BFP}}(x,\xi)\dd\xi \equiv 1$) by compensating the missing boundary weight with a local Dirac delta, $W_\ell^{\mathrm{BFP}}(x,\xi) = \vartheta_\ell(x)\delta(x-\xi) + K_\ell^{\mathrm{unnorm}}(x,\xi)$, as analyzed in detail by Di~Luzio \cite{DiLuzio2005}, Marotti de Sciarra \cite{MarottiDeSciarra2009}, and Ba{\v z}ant, Le, and Hoover \cite{BazantLeHoover2010}.

  \item \textbf{Challamel and Wang on boundary paradoxes of differential nonlocal beams \cite{ChallamelWang2014}:} Challamel and Wang (2008) and Challamel, Zhang, Wang, Reddy, Wang, Michelitsch, and Collet \cite{ChallamelWang2014} demonstrated that Eringen's differential relation $(1 - \ell^2\partial_x^2)M = -EIw''$ produces a non-self-adjoint operator under standard clamped-free boundary conditions ($\int_0^L (M''\delta w - w''\delta M)\dd x = [M'\delta M - M\delta M']_0^L/(EI\ell^{-2}) \neq 0$), and showed that self-adjointness and physical softening require either a discrete-continualized variational boundary condition or a two-phase local/nonlocal mixture.

  \item \textbf{Fern{\'a}ndez-S{\'a}ez, Zaera, Loya, and Reddy \cite{FernandezSaez2016,FernandezSaezZaera2017} and Romano, Barretta, Diaco, and Marotti de Sciarra \cite{Romano2017} on Constitutive Boundary Conditions (CBCs):} Fern{\'a}ndez-S{\'a}ez et al.\ \cite{FernandezSaez2016,FernandezSaezZaera2017} and Romano et al.\ \cite{Romano2017} proved that on a finite interval $[0,L]$, the integral relation $g(x) = \int_0^L K_\ell^{\mathrm{unnorm}}(x,\xi)f(\xi)\dd\xi$ is equivalent to the differential equation $(1-\ell^2\partial_x^2)g = f$ \emph{if and only if} $g$ additionally satisfies the two homogeneous Robin Constitutive Boundary Conditions $(\partial_x \mp \ell^{-1})g = 0$ at $x=0,L$. Without these CBCs, the purely nonlocal ($\gamma = 0$) integral problem is generally over-constrained and admits no smooth solution unless the source satisfies boundary compatibility conditions \cite{Romano2017}.

  \item \textbf{Mikhasev and Nobili \cite{MikhasevNobili2020} and Nobili and Pramanik \cite{NobiliPramanik2025a,NobiliPramanik2025b} on two-phase ($\gamma \in (0,1)$) singular perturbations:} Originally introduced in elasticity by Eringen (1987) and studied by Ba{\v z}ant and Jirásek \cite{BazantJirasek2002}, the two-phase operator $\Lcal_{\gamma,\ell}[f] = \gamma f + (1-\gamma)\Kcal_\ell[f]$ ($\gamma \in (0,1)$) regularizes the Fredholm integral equation of the first kind into one of the second kind. Mikhasev and Nobili \cite{MikhasevNobili2020} proved that as $\gamma \to 0^+$, the fourth-order differential reduction is a singular perturbation whose extra roots scale as $O(\gamma^{-1/2}\ell^{-1})$, spawning exponential boundary layers of thickness $\ell\sqrt{\gamma}$ near $x=0,L$. Nobili and Pramanik \cite{NobiliPramanik2025a,NobiliPramanik2025b} extended two-phase integral elasticity to viscoelasticity and demonstrated that boundary effects under unnormalized kernels mix true physical nonlocality with kernel truncation.

  \item \textbf{Giorgi and Morro \cite{GiorgiMorro2025} on thermodynamic consistency of nonlocal macroscopic models:} Giorgi and Morro \cite{GiorgiMorro2025} systematically analyzed thermodynamic consistency across nonlocal macroscopic theories using the global Clausius--Duhem inequality over $\Omega$, proving that cross-coupling terms between mechanical and thermal fields must arise as conjugate variational derivatives of a common thermodynamic potential.
\end{enumerate}

\subsection{Factual Survey of 29 Spatial-Nonlocal Thermoelasticity Papers}\label{subsec:survey}

To document how these foundational principles are applied in contemporary nonlocal thermoelasticity, we audited a convenience sample of $N = 29$ spatial-nonlocal thermoelasticity papers (published 2016--2026) drawn from a working collection of 152 mechanics PDFs. \textbf{Selection criterion:} Every paper in the 152-PDF collection that solves a spatial-nonlocal thermoelastic problem using Eringen's integral or differential nonlocal operator was included ($N = 29$; 25 solve coupled thermoelastic equations and 4 treat uncoupled thermal loads or uncoupled nonlocal heat conduction; 28 of the 29 are posed on bounded or half-space domains with boundaries). Because this collection was assembled around the author's research area, 13 of the 29 surveyed papers ($44.8\%$) are co-authored by the present author; it is therefore a convenience sample rather than a random sample of the literature.

Table~\ref{tab:census} classifies the 29 surveyed papers according to their constitutive pairing of the stress equation, entropy/energy equation, and boundary closure. All 25 coupled papers in the surveyed sample fall into two differential shortcut classes:

\begin{table}[htbp]
  \centering
  \caption{Classification of the $N = 29$ surveyed spatial-nonlocal thermoelasticity papers (2016--2026) by constitutive pairing and boundary treatment.}\label{tab:census}
  \vspace{4pt}
  \begin{tabular}{p{3.9cm}cccp{6.2cm}}
    \toprule
    \textbf{Constitutive Class} & \textbf{All ($N=29$)} & \textbf{Coupled ($N=25$)} & \textbf{Author Subset} & \textbf{Defining Equations \& Mathematical Status} \\
    \midrule
    \textbf{Class I:} Stress-nonlocal + local constant-strain heating & $12$ ($41.4\%$) & $12$ ($48.0\%$) & $2 / 12$ \cite{BarakGupta2023,GuptaKumar2023a} & $(1-\ell^2\nabla^2)\sigma = E\varepsilon - \beta\theta$ paired with $\rho\eta^{\mathrm{loc}} = \beta\varepsilon + \frac{\rho c_E}{T_0}\theta$. Violates Maxwell symmetry in both Helmholtz and Gibbs forms (Theorem~\ref{thm:no_free_energy}); omits Robin CBCs. \\
    \addlinespace[3pt]
    \textbf{Class II:} All-field identical-kernel cancellation ($\gamma=0$) & $13$ ($44.8\%$) & $13$ ($52.0\%$) & $11 / 13$ \cite{GuptaBarakDas2023b,GuptaBarakDas2024b} & Applies $(1-\ell^2\nabla^2)$ to $\sigma$, $\eta$, and $q$ identically. Cancels out of displacement--temperature PDEs in homogeneous domains; traction BCs become local if applied to $\sigma^{\mathrm{loc}}$. \\
    \addlinespace[3pt]
    \textbf{Class III:} Uncoupled thermal load / nonlocal DPL & $4$ ($13.8\%$) & $0$ ($0.0\%$) & $0 / 4$ & Prescribes $\theta(x,t)$ externally ($T_0\beta\dot{\varepsilon} = 0$) or applies $(1-\ell^2\nabla^2)$ only to uncoupled rigid-conductor heat conduction. \\
    \addlinespace[3pt]
    \textbf{Consistent two-phase (V1 or V2) with CBCs} & $0$ ($0.0\%$) & $0$ ($0.0\%$) & $0$ & Two-phase operator $\Lcal_{\gamma,\ell}$ ($\gamma\in(0,1)$) satisfying $\Lcal_\eta = \Lcal_{\mathrm{th}}^*$ with integral-consistent endpoint CBCs. \\
    \bottomrule
  \end{tabular}
\end{table}

\begin{itemize}[leftmargin=*]
  \item \textbf{Class I (Stress-nonlocal + local heating, 12 of 29 surveyed papers, $41.4\%$ of all, $48.0\%$ of coupled):} Applies $(1-\ell^2\nabla^2)\sigma = E\varepsilon - \beta\theta$ to the stress law while retaining the unmodified local entropy $\rho\eta^{\mathrm{loc}} = \beta\varepsilon + \frac{\rho c_E}{T_0}\theta$ in the energy equation. As shown in Theorem~\ref{thm:no_free_energy}, this hybrid pairing violates Eringen's \cite{Eringen1974} reciprocity identity \eqref{eq:eringen_1974_laws} and admits no free energy in either Helmholtz or Gibbs variables.
  \item \textbf{Class II (All-field identical-kernel cancellation, 13 of 29 surveyed papers, $44.8\%$ of all, $52.0\%$ of coupled):} Applies the identical purely nonlocal ($\gamma = 0$) differential operator $1 - \ell^2\nabla^2$ to the stress, entropy, and heat flux simultaneously. Substituting $(1-\ell^2\nabla^2)\sigma = E\varepsilon - \beta\theta$ into $\partial_x(1-\ell^2\nabla^2)\sigma = (1-\ell^2\nabla^2)(\rho\ddot{u})$ yields $(1-\ell^2\nabla^2)(E\partial_x^2 u - \beta\partial_x\theta - \rho\ddot{u}) = 0$, so the nonlocal factor cancels from the bulk displacement--temperature PDEs unless coupled through nonlocal Cauchy traction BCs and Robin CBCs.
\end{itemize}
We state factually that several of the present author's own earlier papers in the surveyed sample fall into these two classes (specifically, \cite{BarakGupta2023,GuptaKumar2023a} in Class~I and \cite{GuptaBarakDas2023b,GuptaBarakDas2024b} in Class~II). Clarifying how to formulate two-phase nonlocal thermoelasticity without falling into either Class~I or Class~II is the motivation for this note.

\subsection{Scope and Outline of This Clarification Note}\label{subsec:scope}

This note introduces no new physical postulates; it applies the established frameworks of Eringen \cite{Eringen1974}, Polizzotto and Borino--Failla--Parrinello \cite{DiLuzio2005,MarottiDeSciarra2009}, Fern{\'a}ndez-S{\'a}ez et al.\ \cite{FernandezSaez2016}, Romano et al.\ \cite{Romano2017}, Mikhasev and Nobili \cite{MikhasevNobili2020}, Nobili and Pramanik \cite{NobiliPramanik2025b}, and Giorgi and Morro \cite{GiorgiMorro2025} to resolve four specific questions:
\begin{enumerate}[leftmargin=*]
  \item \textbf{Section~\ref{sec:model}:} Formulates the general quadratic nonlocal Helmholtz free energy on $L^2(\Omega)$ and defines the two minimal thermodynamically consistent two-phase ($\gamma \in (0,1)$) realizations: \textbf{Variant~V1} (nonlocal elastic and thermal stress paired with nonlocal dilational entropy) and \textbf{Variant~V2} (nonlocal elastic stiffness paired with local thermal stress and local dilational entropy, as in Pham and Vu \cite{PhamVu2024}).
  \item \textbf{Section~\ref{sec:theorem1}:} Proves \textbf{Theorem~\ref{thm:no_free_energy}}, showing that the Class~I pairing admits no free energy in either Helmholtz $(\varepsilon,\theta)$ or Gibbs $(\sigma,\theta)$ variables, and provides a one-page hand-checkable Maxwell-symmetry and closed-cycle dissipation calculation.
  \item \textbf{Section~\ref{sec:mode_count}:} Proves \textbf{Theorem~\ref{thm:mode_count}}, showing that for $\gamma \in (0,1)$, Variant~V1 is an eighth-order spatial PDE system with $N=4$ wave-mode pairs $\pm k_j(\omega)$ whereas Variant~V2 and Class~I have $N=3$ pairs, and derives the singular-perturbation blow-up $k_4^2(\omega;\gamma) \sim -(\gamma\ell^2)^{-1}$ as $\gamma \to 0^+$.
  \item \textbf{Section~\ref{sec:boundary}:} Derives the exact endpoint Constitutive Boundary Conditions (\textbf{Theorem~\ref{thm:cbc}}) for the unnormalized exponential kernel $K_\ell^{\mathrm{unnorm}}$, the BFP-normalized kernel $W_\ell^{\mathrm{BFP}}$, and the Neumann Green's function $K_\ell^{\mathrm{Neu}}$, and proves (\textbf{Corollary~\ref{cor:artifact}}) that the static $1/\sqrt{\gamma}$ surface thermal-strain ratio between V2 and V1 under uniform heating is a kernel-truncation artifact of $K_\ell^{\mathrm{unnorm}}$ that vanishes identically under $W_\ell^{\mathrm{BFP}}$ and $K_\ell^{\mathrm{Neu}}$.
  \item \textbf{Section~\ref{sec:calibration}:} Evaluates the bulk attenuation differences for fcc Aluminum using the atomistic 2NN MEAM characteristic length $\bar{\ell}_1 = 0.1203\;\mathrm{nm}$ of Lazar, Agiasofitou, and B{\"o}hlke \cite{Lazar2022}, demonstrating plainly that within the frequency regime where Fourier heat conduction is physically valid, relative differences between V1, V2, and Class~I are $O(10^{-4})$ or smaller. Section~\ref{sec:limitations} summarizes the limitations and modeling guidelines, and Appendix~\ref{app:proofs} collects the full mathematical proofs.
\end{enumerate}
